\section*{Chapitre \ref{ch:qm:state-space-models-and-the-kalman-filter} --- Modèles espace-état et filtre de Kalman}

\begin{solution}{exo:qm:state-space-models-and-the-kalman-filter:1}
$q = 0.04$~: $p = (0.04 + \sqrt{0.0016 + 0.16})/2 = 0.221$ et $K = 0.221/1.221 = 0.181$. Les poids des anciennes observations décroissent comme $0.819^k$, ce qui divise par deux en
$\ln\frac12/\ln0.819 = 3.5$ jours.
\end{solution}

\begin{solution}{exo:qm:state-space-models-and-the-kalman-filter:2}
Variance prédite $0.01 + 0.0002 = 0.0102$. \omterm{def:qm:state-space-models-and-the-kalman-filter:kf}{Innovation} $v = 1.2 - 0.8 \times 2 = -0.4$, $F = 4 \times 0.0102 + 0.5 = 0.5408$, \omterm{def:qm:state-space-models-and-the-kalman-filter:kf}{gain de Kalman} $K = 0.0102 \times 2/0.5408 =
0.0377$. Ratio mis à jour $0.8 - 0.0377 \times 0.4 = 0.785$, variance $0.0102(1 - 0.0377 \times 2) = 0.0094$.
\end{solution}

\begin{solution}{exo:qm:state-space-models-and-the-kalman-filter:3}
Sans observation, il n'y a pas de mise à jour, seulement l'étape de prédiction, qui ajoute le bruit d'état~: dans le \omterm{def:qm:state-space-models-and-the-kalman-filter:ssm}{modèle à niveau local}, la variance augmente de $Q$ ce jour-là,
au lieu de se réduire d'abord avec la mise à jour. L'incertitude traduit le fait que l'état a continué à bouger alors que rien n'était observé.
\end{solution}

\begin{solution}{exo:qm:state-space-models-and-the-kalman-filter:4}
Pour $q$ petit, $p = \frac12(q + \sqrt{q^2 + 4q}) \approx \sqrt q$ et $K = p/(1 + p) \approx \sqrt q$. Le niveau filtré est $a_{t+1} = (1 - K)a_t + Ky_t$, une moyenne mobile exponentielle de paramètre $\lambda
= 1 - K$~; pour $q = 0.0004$, $K = 0.0198$ et $\lambda = 0.980$.
\end{solution}

\begin{solution}{exo:qm:state-space-models-and-the-kalman-filter:5}
État $\alpha_t = x_t$, $Z = 1$, $T = \phi$, $Q = \Var(\eta)$, $H = \Var(\varepsilon)$. La variance prédite en régime permanent vérifie $P = \phi^2\bigl(P - P^2/(P + H)\bigr) + Q$.
\end{solution}

\begin{solution}{exo:qm:state-space-models-and-the-kalman-filter:6}
$1/(0.25 + 0.09 + 0.01 + 0.0025 + 0.0025) = 1/0.355 = 2.8$ particules équivalentes sur cinq.
\end{solution}

\begin{solution}{exo:qm:state-space-models-and-the-kalman-filter:7}
\texttt{tradeoff()}~: $Q \times 0.1$~: 0.987 de l'erreur de la fenêtre, demi-réponse de 96 jours~; $\times 1$~: 0.993, 30 jours~; $\times 10$~: 1.006, 9.6 jours~; $\times 100$~: 1.062, 3.0 jours.
Un choix plus lent que celui de la vraisemblance est ici légèrement meilleur~; un choix plus rapide coûte du bruit.
\end{solution}

\begin{solution}{exo:qm:state-space-models-and-the-kalman-filter:8}
Le ratio lissé pour le jour $t$ utilise les cours postérieurs au jour $t$, si bien que le backtest se couvre avec une information dont le desk n'aurait pas pu disposer. Utiliser le ratio filtré à partir de la
clôture précédente~; le lissage relève de la recherche sur ce qu'était la relation, et non d'une simulation de trading.
\end{solution}

\begin{solution}{pb:qm:state-space-models-and-the-kalman-filter:1}
\textbf{1.} Le cours spot du WTI publié par l'EIA s'est établi à $-36.98$~dollars~; le WTI a varié de $-55.29$ et $+45.89$ les 20 et 21 avril, tandis que le Brent a varié de $-2.39$ et $-8.24$.
Ces deux jours sont traités comme des données manquantes~: le filtre ignore la mise à jour, et la régression glissante ainsi que l'évaluation écartent ces jours.
\textbf{2.} 0.41 en 2014, 0.94 en 2021 et 1.01 en 2026.
\textbf{3.} La volatilité quotidienne du WTI va de 0.78~dollars (2017) à 3.62 (2026)~; avec $H$ constant, les périodes volatiles dominent la vraisemblance et tirent $Q$ à la hausse~;
la normalisation rend le bruit d'observation plus proche de l'homoscédasticité.
\textbf{4.} 3.324 sans couverture et 1.329 avec le ratio des 60 premiers jours maintenu fixe (dollars au carré par baril et par jour, 4\,035 jours).
\textbf{5.} 1.276.
\textbf{6.} $H = 0.508$, $Q = 0.000238$, \omterm{def:qm:state-space-models-and-the-kalman-filter:snr}{rapport signal sur bruit} 0.00047.
\textbf{7.} $\ln2/\sqrt{0.00047 \times 1.12} = 30$ jours, soit autant que les 30 de la fenêtre de 60 jours.
\textbf{8.} 1.267, soit une réduction de 0.7~\%.
\textbf{9.} Le $Q$ estimé est dix fois plus grand (0.0026 contre $H = 1.14$ en unités de dollars) et le filtre couvre 1.4~\% moins bien que la fenêtre.
\textbf{10.} À environ $\pm 0.2$~: l'écart type du ratio filtré vaut environ 0.10.
\textbf{11.} $\times 10$~: 0.6~\% d'erreur en plus que la fenêtre, demi-réponse de 9.6 jours~; $\times 100$~: 6.2~\% en plus, 3.0 jours.
\textbf{12.} Environ cinquante jours pour le filtre de vraisemblance (50) et pour la fenêtre (49)~; huit jours avec un $Q$ cent fois plus grand (18 avec dix fois plus).
\textbf{13.} Depuis $H = 1$, $Q = 0.01$, l'algorithme atteint $-4\,524$ après dix pas et $-4\,502$ après quarante, contre le maximum $-4\,470$~: la vraisemblance est plate le long d'une crête
où $Q$ et $H$ se compensent, ce qui rend les pas de l'EM courts.
\textbf{14.} Filtré 1.20, fenêtre glissante 1.28 (22 septembre 2026).
\textbf{15.} Quand les déplacements sont grands et rares par rapport au bruit quotidien (une rupture structurelle), ce qu'une marche aléatoire à pas gaussiens ne décrit pas~; un modèle avec
des sauts occasionnels du ratio ou un détecteur de \omterm{def:qm:robust-statistics-and-heavy-tails:influence}{points de rupture} réagirait rapidement, seulement lorsque c'est nécessaire.
\textbf{16.} Seulement pour la bande d'incertitude, le réglage rigoureux et le traitement des jours manquants~; pas pour l'erreur de couverture, qui est presque identique.
\textbf{17.} Synchroniser les cours (les deux évaluations spot sont relevées à des moments différents), ajouter la variation retardée du WTI, et se couvrir avec des contrats à terme plutôt
qu'avec les cours spot.
\textbf{18.} Parce qu'il utilise des données futures~: tout backtest qui s'en sert se couvre avec une connaissance a posteriori.
\textbf{19.} Résultat retenu~: \emph{le ratio de couverture qui a bougé}~: sur les données 2010--2026, le ratio de couverture de Kalman estimé par maximum de vraisemblance réduit de
0.7~\% la variance de l'erreur de couverture par rapport à la régression glissante sur 60 jours, avec un \omterm{def:qm:state-space-models-and-the-kalman-filter:snr}{rapport signal sur bruit} estimé à 0.00047, ce qui rend le filtre aussi lent que la fenêtre.
\textbf{20.} Le \omterm{def:qm:state-space-models-and-the-kalman-filter:snr}{rapport signal sur bruit}, soit le rapport entre le bruit d'état et le bruit d'observation, que la vraisemblance estime à partir des données.
\end{solution}

\begin{solution}{iq:qm:state-space-models-and-the-kalman-filter:1}
On croit qu'une grandeur qu'on ne voit pas, comme un ratio de couverture, dérive lentement~; chaque jour, on la prédit, on mesure à quel point la prédiction des cours du jour était fausse, et on déplace sa croyance
vers les observations, d'une fraction qui dépend du bruit des cours comparé à la vitesse à laquelle la grandeur cachée se déplace. On suit aussi l'incertitude de cette croyance.
C'est la façon optimale de procéder lorsque tout est linéaire et gaussien.
\iqlookfor{Prédire, comparer, corriger en proportion du bruit relatif, et suivre l'incertitude.}
\end{solution}

\begin{solution}{iq:qm:state-space-models-and-the-kalman-filter:2}
A priori $\alpha \sim \mathcal N(a, P)$, observation $y = z\alpha + \varepsilon$, $\varepsilon \sim \mathcal N(0, H)$. $(\alpha, y)$ sont conjointement gaussiens, avec $\Cov = Pz$ et $\Var(y) = z^2P + H$.
Par conditionnement~: $a' = a + \frac{Pz}{z^2P + H}(y - za)$ et $P' = P - \frac{P^2z^2}{z^2P + H}$~; le \omterm{def:qm:state-space-models-and-the-kalman-filter:kf}{gain de Kalman} vaut $K = Pz/(z^2P + H)$.
\iqlookfor{L'argument de normalité conjointe, et non une formule apprise par cœur.}
\end{solution}

\begin{solution}{iq:qm:state-space-models-and-the-kalman-filter:3}
L'estimer, avec le bruit d'observation, par maximum de vraisemblance à partir de la \omterm{def:qm:state-space-models-and-the-kalman-filter:ped}{décomposition par erreurs de prédiction}, puis vérifier la couverture hors échantillon et le compromis vitesse--bruit~;
le fixer à la main pour «~réagir en quelques jours~» revient à acheter du bruit. Normaliser d'abord le bruit d'observation par la volatilité, faute de quoi l'ajustement est dominé par les périodes volatiles.
\iqlookfor{Réglage par vraisemblance, compromis, et hétéroscédasticité.}
\end{solution}

\begin{solution}{iq:qm:state-space-models-and-the-kalman-filter:4}
Le filtrage conditionne sur les données jusqu'à $t$~; le lissage, sur l'ensemble des données, passées et futures. La trajectoire lissée est meilleure pour décrire l'histoire et fausse pour les simulations de trading,
où elle constitue un biais d'anticipation.
\iqlookfor{Ensembles d'information et biais d'anticipation.}
\end{solution}

\begin{solution}{iq:qm:state-space-models-and-the-kalman-filter:5}
Lorsque le modèle est non linéaire ou non gaussien~: volatilité stochastique, sauts, bruit à queue épaisse, régimes discrets. Il peut dégénérer (quelques particules portent le
poids) lors d'observations surprenantes, sa vraisemblance est bruitée (ce qui complique l'estimation des paramètres), et son coût croît avec la dimension de l'état.
\iqlookfor{La non-linéarité comme raison, la dégénérescence et le bruit de Monte-Carlo comme risques, la \omterm{def:qm:bayesian-methods:ess}{taille d'échantillon effective} comme indicateur de suivi.}
\end{solution}

\begin{solution}{iq:qm:state-space-models-and-the-kalman-filter:6}
Pour toute loi $q$ des variables cachées, $\ln p(y \mid \theta) = \E_q[\ln p(y, \alpha \mid \theta)] - \E_q[\ln q] + \mathrm{KL}(q\,\|\,p(\alpha \mid y, \theta))$. L'étape E fixe $q$ à la
\omterm{def:qm:bayesian-methods:prior}{loi a posteriori} en $\theta_k$, ce qui annule le terme de KL~; la borne touche donc la vraisemblance en ce point~; l'étape M fait croître la borne~; la vraisemblance en $\theta_{k+1}$ est au moins
la borne relevée, donc au moins la vraisemblance en $\theta_k$.
\iqlookfor{La borne inférieure et l'écart de KL.}
\end{solution}
